\section{Discussion and Conclusion}\label{sec:discussion}
% Thai version: 05_discussion_th.tex (main_short_th.tex); keep the two in step.
% Replaces the former 06_discussion.tex + 08_limitations.tex + 09_conclusion.tex of the short version.
% Every number here is reported in Section IV; this chapter interprets, relates to prior work and qualifies.

This section relates the results to prior work, interprets the error that remains, draws implications, states the
limitations and future work, and concludes objective by objective.

\subsection{Relation to prior work}\label{sec:prior}
\mbox{DA-v2} Metric-Indoor learned metres from the rendered rooms of Hypersim~\cite{da2,hypersim} and placed every real
test room 1.18 to 1.45 times too far. ZoeDepth~\cite{zoedepth} and \mbox{DA-v2} adapt their metric heads to a target
domain by fine-tuning on labelled data; our result shows that the same remedy works with the target device's own LiDAR
as the labels. As the DINOv2 authors did for depth~\cite{dinov2}, we trained only the decoder on the frozen encoder,
and the shape of the prediction barely changed (per-frame oracle AbsRel 0.035 against 0.042), which fits the view that
the encoder's features transfer~\cite{yosinski} and the head mainly learned the scale.

Sensor depth has long served as a training label, from Kinect indoors to laser scanners on
cars~\cite{nyuv2,kitti,eigen14}. The roughly \SI{1}{cm} that phone LiDAR adds to the Faro depth in our pipeline agrees
with the centimetre-level accuracy reported against surveying methods~\cite{luetzenburg,spreafico}, and its error as a
label (AbsRel \rsLIDARABSREL{}) is an order of magnitude below the error left in the fine-tuned model
(\rsFTLIDARALLABSREL{}), which may explain why a laser teacher gave no measurable gain. Prompt Depth
Anything~\cite{promptda} and the upsampling task of ARKitScenes~\cite{arkitscenes} use the LiDAR at run time; our model
needs only a colour image, at a large price (\SI{\rsFTLIDARALLCM}{cm} against \SI{\rsLIDARCM}{cm}).

Camera-aware models~\cite{metric3d,unidepth,depthpro} treat the unknown focal length as the main obstacle to metric
depth. For Depth Pro on these $640\times480$ indoor frames, the true focal length did not make it usable, and its error
without scale stayed several times that of \mbox{DA-v2} Metric-Indoor; we did not test Metric3D or UniDepth, so this
does not extend to the approach in general. Per-image benchmarks cannot show whether the scale stays the same across
the frames of a scan (Section~\ref{sec:related}); we found that it does not, for relative and metric models alike.
Video Depth Anything~\cite{videoda} targets such consistency across video frames, and sparse metric points, used as in
CNN-SLAM and MonoSDF~\cite{cnnslam,monosdf}, came within \SI{\rsSPARSEGAP}{cm} of the per-frame oracle in their
upper-bound form. Learned reconstruction systems~\cite{atlas,neuralrecon,simplerecon} use several frames per prediction
and other data, so our numbers are not comparable with theirs.

\subsection{Why a scale error remains}\label{sec:whyft}
Fine-tuning learned the domain's multiplier: \ftmain{} without calibration (AbsRel 0.134) is on a par with the
pretrained model given one scale and offset per room fitted on ground truth (0.127), so what a one-off calibration
provides has moved into the weights of the head. The swing from frame to frame did not shrink (1.52 against 1.55). A
likely reason follows from the pinhole model (Section~\ref{sec:related}): views such as a close-up of a table top or a
blank wall look alike at different distances and give no cue to absolute size, so the model falls back on a typical
indoor distance; the negative correlation between the needed scale and the depth of the view, found in every room, fits
this reading. If it is right, more or better labels will not remove the swing, and metric data at run time will be
needed. TSDF fusion~\cite{curless96} then averages the same surface placed at several distances into a thicker, blurred
surface, consistent with a moderate per-frame error becoming \SI{\rsFTLIDARALLCM}{cm} in Chamfer distance.

\subsection{Implications for practice}\label{sec:business}
For an app on a device without LiDAR, the results suggest an order of work: per-frame metric scale first (within
\SI{\rsSPARSEGAP}{cm} of the per-frame oracle, against \SI{\rsSCENECM}{cm} with one calibration per room, an upper bound
because the points came from LiDAR, not from VIO), then a fine-tuned model, which removes the domain bias with nothing at
run time, and model size last (\mbox{DA-v2} Small was \SI{1.0}{cm} worse than Large under the per-frame oracle and 6.4
times faster). The labels can be collected from Pro devices during normal use, since ARKit delivers colour, LiDAR depth
and pose together. At 3.6 times the video length on a laptop, the pipeline as tested runs offline.

\subsection{Limitations}\label{sec:limitations}
The evaluation is small: six test rooms in 3D, one used during development, and 26 rooms per frame, all homes captured
with one iPad model; offices, empty rooms, low light and other phone cameras are absent, and the bathroom with a mirror
was the hardest room for every source with the right scale (Table~\ref{tab:perroom}). With six rooms the Wilcoxon test
detects only a difference present in every room, so the teacher comparison cannot rule out differences of the size seen
in the means, one to two centimetres. The main comparison was fixed after the six test rooms had been evaluated; the
\rsVN{} rooms of Exp.~7 were added afterwards so that the conclusion also rests on rooms that decided nothing. The
reference shares the ARKit poses of every row, so a real phone system would add pose drift, and completeness is relative
to what Faro saw. Only the head of one model was fine-tuned, with one recipe, and \ftfaro{} and \ftlidar{} differ in
the number of views as well as the labels. Checkpoints were chosen with Faro depth, also for the LiDAR-taught models, on
three rooms; \ftmain{} became the main model because it was best there and its labels need no laser scanner, but that
ranking did not reproduce on the test rooms, and its validation AbsRel changed by up to 0.1 between checkpoints 500
steps apart. The once-per-room and sparse-point
alignments are upper bounds, calibrated on ground truth and on clean LiDAR samples. Depth Pro was run through one
implementation, and its focal-length check covered five frames of one room. The thresholds that link recall to use
cases are our illustration, not an industry standard, and times are unoptimised and meaningful only as ratios.

\subsection{Future work}\label{sec:future}
In order of expected impact, future work should (1)~combine \ftmain{} with per-frame metric points from the device's
tracker, for which Section~\ref{sec:scale} gives AbsRel 0.051 with the right scale on every frame; (2)~choose
checkpoints with LiDAR depth, so that no step needs a laser scanner, and compare teachers on identical frames and more
test rooms; (3)~train and evaluate on more phone cameras, including Android; (4)~evaluate with real ARKit or ARCore
feature points and a reference captured by the app itself; and (5)~test fine-tuning more of the network, and video
models that keep the scale consistent across frames~\cite{videoda}.

\subsection{Conclusion}\label{sec:conclusion}
We asked whether fine-tuning a metric monocular depth model with a phone's own LiDAR makes it more accurate than its
pretrained self against a laser scanner. For objective~(1) it did, on the rooms tested and by the criterion of the
Introduction: \ftmain{} lowered the Chamfer distance from \rsMETRICCM{} to \SI{\rsFTLIDARALLCM}{cm} in all six test rooms
and AbsRel from \rsVMETRICABSREL{} to \rsVFTALLABSREL{} in \rsVWINS{} of \rsVN{} further rooms, with no calibration or
sensor at run time and with weights learned from LiDAR alone, although \wone{} improved in only five of the six test
rooms. For objective~(2), laser labels were not measurably better than LiDAR labels, nor 24 rooms better than 10, within
the limited power of six rooms, and the true focal length did not make Depth Pro competitive. For objective~(3), the
model lies between a relative model calibrated once per room and one aligned on sparse points in every frame, serves
none of the three use cases, and is no substitute for iPad LiDAR (\SI{\rsLIDARCM}{cm}). For objective~(4), its remaining
error is a scale that varies with the view rather than drifting with time, whose bias fine-tuning removed but whose
spread it did not. Per-frame metric points from the device's tracker are the most promising next step; their
combination with the fine-tuned model remains to be tested.
